\documentclass[11pt,a4paper]{article}
\usepackage[margin=2.5cm]{geometry}
\usepackage{amsmath,amssymb,booktabs,xcolor,enumitem}
\usepackage[normalem]{ulem}
\usepackage[hidelinks]{hyperref}
\setlength{\parskip}{0.6em}\setlength{\parindent}{0pt}
\definecolor{newc}{HTML}{1f6f43}
\newcommand{\new}[1]{\textcolor{newc}{#1}}
\newcommand{\old}[1]{\textcolor{gray}{\sout{#1}}}
\newcommand{\oldmath}[1]{\textcolor{gray}{#1}}
\newcommand{\mockup}[1]{\textit{#1}}
\title{Proposed change [A93], the arrivals of Appendix C follow the refill rule of the code}
\date{13 November 2026, version 1}
\begin{document}
\maketitle

\section*{How to read this}
Part~1 says what is wrong. Part~2 shows Appendix~C's ``The arrivals'' as it would read, and the
notation rows. \new{Green} is new text, \old{struck} text goes, a removed equation is grey and marked
``removed''. Part~3 has the questions. Nothing is changed in the thesis yet.

\section*{Part 1. What is wrong}

Appendix~C defines how much arrives by the \emph{top-up} rule. Each SKU is filled up to
$(1+\varepsilon)$ times last period's orders, never below a floor $x_{\min}$. The code has that rule
(\texttt{stock\_rule: topup}), but the thesis runs use the other one, \texttt{rule: refill} with
\texttt{refill\_units: 10} in every \texttt{configs/step12/fleet\_<m>.yaml}.

Under the refill rule each SKU counts the units ordered since its last refill, from the moment a line is
released (\texttt{multi\_agent\_env.py}, \texttt{\_ordered\_since\_refill}). At an arrival it gets one
refill of $Q$ units for every $Q$ it has counted, and the rest carries over
(\texttt{storage/arrival.py}, \texttt{refill\_amounts}). Every refill is one put-away to one shelf.
Counting the units on their way and leaving out the units lines wait for, a SKU's stock then stays
between its starting stock less $Q$ and its starting stock.

Appendix~G (``Setting'') and the Problem Formulation already describe refills. Only Appendix~C and
two rows of the notation table still describe the top-up rule. $\hat\rho^{\mathrm{unit}}_t$ is used
only there, so its row goes too.

\section*{Part 2. The text}

\subsection*{Appendix C, ``The arrivals''}

How much arrives is set by the orders alone, so it is the same under every
rule. \new{Stock comes back a refill at a time. Each SKU counts the units
ordered since its last refill, from the moment a line is released, whether
the units are picked yet or still wait. At a storage update a SKU gets one
refill of $Q$ units for every $Q$ it has counted, and what is left over
counts towards the next one. So with $Q = 10$, if 23 units of tea were
ordered since the last refill, two refills of tea arrive and 3 units carry
over. A SKU's stock, counting the units on their way and leaving out the
units lines wait for, then stays between its starting stock less $Q$ and
its starting stock. This is the $(R, nQ)$ policy of inventory control
(Axs\"ater, 2015, \texttt{axsater2015inventory}), with the reorder point $R$ one refill below
the starting stock.}
\old{Each SKU is filled up to what the last period's orders suggest, with a
margin, and never to fewer than a floor of units. The units that lines still
wait for are added on top. Arrivals count units, while $\hat\rho_t$ counts
lines, so a second estimate counts the units ordered. So if 30 units of tea
were ordered in the last period, the margin is $1$ and the floor is $8$, tea
is filled up to $60$ units plus the units still waited for. Units still on
their way to a shelf count as stock, so stock already ordered is not ordered
again, and the arrivals stay the same under every rule.}

\new{Formally, write $y^{\mathrm{ord}}_t(k)$ for the units of SKU $k$ ordered
since its last refill, just before the update at $t$, with ``ord'' for
ordered. The \emph{arrivals} are}
\old{Formally, with margin $\varepsilon \ge 0$, floor $x_{\min}$,
$\hat\rho^{\mathrm{unit}}_t(k)$ the units of $k$ ordered per step in the
window, and $q^{\mathrm{wait}}_t(k)$ the units of $k$ that lines wait for at
$t$, the \emph{arrivals} are}
\begin{equation*}
  \new{y_t(k) = Q \left\lfloor \frac{y^{\mathrm{ord}}_t(k)}{Q} \right\rfloor ,}
  \tag{\texttt{eq:arrivals}}
\end{equation*}
\new{and $y^{\mathrm{ord}}_t(k) - y_t(k)$, fewer than $Q$ units, carries over.}

\oldmath{Removed:
$y_t(k) = \max\bigl(0,\; x^{\mathrm{up}}_t(k) + q^{\mathrm{wait}}_t(k)
- \sum_{v} (x_t^{-}(k,v) - r_t(k,v) + x^{\mathrm{in},-}_t(k,v))\bigr)$,
$x^{\mathrm{up}}_t(k) = \max(x_{\min},\, \lceil (1+\varepsilon)\, \Delta\,
\hat\rho^{\mathrm{unit}}_t(k) \rceil)$,} \old{where $\Delta\,\hat\rho^{\mathrm{unit}}_t(k)$ is the
number of units of $k$ ordered in the period.}

\subsection*{Notation table (Appendix A)}
\begin{center}
\begin{tabular}{lll}
\toprule
symbol & name & meaning \\
\midrule
\old{$\varepsilon \ge 0$} & \old{stock margin} & \old{Stock kept above the last period's orders} \\
\old{$x_{\min}$} & \old{stock floor} & \old{Fewest units a SKU is filled up to} \\
\old{$\hat\rho^{\mathrm{unit}}_t(k)$} & \old{unit demand estimate} & \old{Units of SKU $k$ ordered per step \ldots} \\
\new{$y^{\mathrm{ord}}_t(k)$} & \new{units ordered} & \new{Units of SKU $k$ ordered since its last refill, before the update at $t$} \\
\bottomrule
\end{tabular}
\end{center}
\mockup{The row for $Q$, the refill size, already exists. The same rows change in the unused copy
\texttt{Chapters/2.Notation.tex}, which follows Appendix~A.}

\section*{Part 3. Questions}
\begin{enumerate}[leftmargin=1.4em]
\item \textbf{The citation.} The code calls the rule the $(R, nQ)$ policy. I could not open Axs\"ater's
  book to check its definition there. Keep the citation and the clause on the reorder point if it
  matches, or say only ``ordering in fixed batches of $Q$'' without a source?
\item \textbf{The worked numbers.} 23 units ordered, two refills, 3 carry over. Keep the example, or
  leave it to the equation?
\item \textbf{The stock bound.} The sentence that a SKU's stock stays between its starting stock less
  $Q$ and its starting stock comes from the code's comment. It explains why a home shelf always has
  room (Appendix~G, ``Shelf capacity''). Keep it here, or only in Appendix~G?
\end{enumerate}

\end{document}
